% GENERATED FILE. Edit canonical lessons, then run npm run generate:books.
% canonical-source-hash: fnv1a64:51410968d6a72bdc
% canonical-lessons: ML-C520-anukampa, ML-C520-sahatapam, ML-C520-vinnuka, ML-C520-asvasthata, ML-C520-vitumpuka

\chapter{Compassion and Panic}
\label{ch:ml-compassion-and-panic}
\hlchaptermodality{520}

\begin{chapteropening}
\textbf{By the end of this chapter:} \emph{I can say compassion, sympathy, to throb, unease and to sob.}

Complete the last lesson of chapter 520: I can say compassion, sympathy, to throb, unease and to sob.
\end{chapteropening}

\section[\texorpdfstring{\ml{അനുകമ്പ} (anukampa)}{അനുകമ്പ (anukampa)}]{\ml{അനുകമ്പ} (anukampa) — compassion}

\label{lesson:ML-C520-anukampa}



\begin{quote}
Before the new one: say the Malayalam for tender affection, then the Malayalam for to grumble.
\end{quote}

\subsection*{You'll want to know: \ml{അനുകമ്പ}}
\textbf{\ml{അനുകമ്പ}} — \emph{anukampa} — \textquotedblleft{}compassion\textquotedblright{}.

\begin{quote}
\ml{രോഗികളോട്} \ml{ആ} \ml{നഴ്സിന്} \ml{വലിയ} \ml{അനുകമ്പ} \ml{ഉണ്ടായിരുന്നു}.
\end{quote}

\begin{quote}
\emph{rōgikaḷōṭŭ ā naḻsinŭ valiya anukampa uṇṭāyirunnu.}
\end{quote}

\begin{quote}
That nurse had great compassion for the patients.
\end{quote}

A word of writing and speeches; the one pitied takes -\ml{ഓട്}.

\begin{cousinweb}
\textbf{\ml{അനുകമ്പ}} (\emph{anukampa}) is borrowed from Sanskrit \emph{anukampā}, \textquotedblleft{}sympathy, compassion\textquotedblright{} (Monier-Williams), from \emph{anu-kamp}, \textquotedblleft{}to tremble along with someone\textquotedblright{}.
\end{cousinweb}

\begin{grammarlens}[title={What you've built}]
One more word for describing what you lived and felt.
\end{grammarlens}

\subsection*{Your turn}
\begin{itemize}
\raggedright
  \item \emph{Say it:} \emph{anukampa}
  \item \emph{Say it:} \emph{anukampa}, once more
  \item \emph{Say it:} \emph{rōgikaḷōṭŭ ā naḻsinŭ valiya anukampa uṇṭāyirunnu.} — That nurse had great compassion for the patients.
  \item \emph{Recall:} say the Malayalam for tender affection, then the Malayalam for to grumble, then say \emph{anukampa} again
\end{itemize}

\begin{tcolorbox}[breakable,colback=teal!4,colframe=teal!35!black,title={Before you move on}]
What does \ml{അനുകമ്പ} mean? (\textquotedblleft{}Compassion\textquotedblright{}.) Say it once more.
\end{tcolorbox}
\section[\texorpdfstring{\ml{സഹതാപം} (sahatāpaṁ)}{സഹതാപം (sahatāpaṁ)}]{\ml{സഹതാപം} (sahatāpaṁ) — sympathy; pity}

\label{lesson:ML-C520-sahatapam}



\begin{quote}
Before the new one: say the Malayalam for to grumble, then the Malayalam for compassion.
\end{quote}

\subsection*{You'll want to know: \ml{സഹതാപം}}
\textbf{\ml{സഹതാപം}} — \emph{sahatāpaṁ} — \textquotedblleft{}sympathy; pity\textquotedblright{}.

\begin{quote}
\ml{അയൽക്കാർ} \ml{ഞങ്ങളോട്} \ml{സഹതാപം} \ml{കാണിച്ചു}.
\end{quote}

\begin{quote}
\emph{ayalkkār ñaṅṅaḷōṭŭ sahatāpaṁ kāṇiccu.}
\end{quote}

\begin{quote}
The neighbours showed us sympathy.
\end{quote}

\ml{സഹ}- with + \ml{താപം} heat, pain: feeling the pain with someone.

\begin{cousinweb}
\textbf{\ml{സഹതാപം}} (\emph{sahatāpaṁ}) is a learned compound built from Sanskrit \emph{saha}, \textquotedblleft{}with\textquotedblright{}, and \emph{tāpa}, \textquotedblleft{}heat, pain\textquotedblright{} (Monier-Williams).
\end{cousinweb}

\begin{grammarlens}[title={What you've built}]
One more word for describing what you lived and felt.
\end{grammarlens}

\subsection*{Your turn}
\begin{itemize}
\raggedright
  \item \emph{Say it:} \emph{sahatāpaṁ}
  \item \emph{Say it:} \emph{sahatāpaṁ}, once more
  \item \emph{Say it:} \emph{ayalkkār ñaṅṅaḷōṭŭ sahatāpaṁ kāṇiccu.} — The neighbours showed us sympathy.
  \item \emph{Recall:} say the Malayalam for to grumble, then the Malayalam for compassion, then say \emph{sahatāpaṁ} again
\end{itemize}

\begin{tcolorbox}[breakable,colback=teal!4,colframe=teal!35!black,title={Before you move on}]
What does \ml{സഹതാപം} mean? (\textquotedblleft{}Sympathy; pity\textquotedblright{}.) Say it once more.
\end{tcolorbox}
\section[\texorpdfstring{\ml{വിങ്ങുക} (viṅṅuka)}{വിങ്ങുക (viṅṅuka)}]{\ml{വിങ്ങുക} (viṅṅuka) — to throb; to swell with grief}

\label{lesson:ML-C520-vinnuka}



\begin{quote}
Before the new one: say the Malayalam for compassion, then the Malayalam for sympathy.
\end{quote}

\subsection*{You'll want to know: \ml{വിങ്ങുക}}
\textbf{\ml{വിങ്ങുക}} — \emph{viṅṅuka} — \textquotedblleft{}to throb; to swell with grief\textquotedblright{}.

\begin{quote}
\ml{സുഹൃത്ത്} \ml{പോയപ്പോൾ} \ml{അവളുടെ} \ml{മനസ്സ്} \ml{വിങ്ങി}.
\end{quote}

\begin{quote}
\emph{suhṛttŭ pōyappōḷ avaḷuṭe manassŭ viṅṅi.}
\end{quote}

\begin{quote}
When her friend left, her heart ached.
\end{quote}

Past: \ml{വിങ്ങി}, \emph{viṅṅi}.

Of a wound it means to throb and swell; of the heart, to ache with held-back grief.

\begin{cousinweb}
\textbf{\ml{വിങ്ങുക}} (\emph{viṅṅuka}) is native Dravidian and seems to be a short-vowel form of \ml{വീങ്ങുക}, \textquotedblleft{}to swell\textquotedblright{}: DEDR 5448 lists that verb with \ml{വീക്കം}, a swelling, and the throbbing of a wound (DEDR 5448).
\end{cousinweb}

\begin{grammarlens}[title={What you've built}]
One more word for describing what you lived and felt.
\end{grammarlens}

\subsection*{Your turn}
\begin{itemize}
\raggedright
  \item \emph{Say it:} \emph{viṅṅuka}
  \item \emph{Say it:} \emph{viṅṅuka}, once more
  \item \emph{Say it:} \emph{suhṛttŭ pōyappōḷ avaḷuṭe manassŭ viṅṅi.} — When her friend left, her heart ached.
  \item \emph{Recall:} say the Malayalam for compassion, then the Malayalam for sympathy, then say \emph{viṅṅuka} again
\end{itemize}

\begin{tcolorbox}[breakable,colback=teal!4,colframe=teal!35!black,title={Before you move on}]
What does \ml{വിങ്ങുക} mean? (\textquotedblleft{}To throb; to swell with grief\textquotedblright{}.) Say it once more.
\end{tcolorbox}
\section[\texorpdfstring{\ml{അസ്വസ്ഥത} (asvasthata)}{അസ്വസ്ഥത (asvasthata)}]{\ml{അസ്വസ്ഥത} (asvasthata) — unease; restlessness}

\label{lesson:ML-C520-asvasthata}



\begin{quote}
Before the new one: say the Malayalam for sympathy, then the Malayalam for to throb.
\end{quote}

\subsection*{You'll want to know: \ml{അസ്വസ്ഥത}}
\textbf{\ml{അസ്വസ്ഥത}} — \emph{asvasthata} — \textquotedblleft{}unease; restlessness\textquotedblright{}.

\begin{quote}
\ml{ചൂട്} \ml{കാരണം} \ml{രാത്രി} \ml{മുഴുവൻ} \ml{എനിക്ക്} \ml{അസ്വസ്ഥത} \ml{തോന്നി}.
\end{quote}

\begin{quote}
\emph{cūṭŭ kāraṇaṁ rātri muḻuvan enikkŭ asvasthata tōnni.}
\end{quote}

\begin{quote}
Because of the heat I felt restless all night.
\end{quote}

From \ml{അസ്വസ്ഥം}, not at ease, with -\ml{ത}; \ml{അ}- turns a Sanskrit word into its opposite.

\begin{cousinweb}
\textbf{\ml{അസ്വസ്ഥത}} (\emph{asvasthata}) is borrowed from Sanskrit \emph{asvasthatā}, \textquotedblleft{}illness\textquotedblright{} (Monier-Williams), the opposite of \emph{svastha}, \textquotedblleft{}being in oneself, well\textquotedblright{}.
\end{cousinweb}

\begin{grammarlens}[title={What you've built}]
One more word for describing what you lived and felt.
\end{grammarlens}

\subsection*{Your turn}
\begin{itemize}
\raggedright
  \item \emph{Say it:} \emph{asvasthata}
  \item \emph{Say it:} \emph{asvasthata}, once more
  \item \emph{Say it:} \emph{cūṭŭ kāraṇaṁ rātri muḻuvan enikkŭ asvasthata tōnni.} — Because of the heat I felt restless all night.
  \item \emph{Recall:} say the Malayalam for sympathy, then the Malayalam for to throb, then say \emph{asvasthata} again
\end{itemize}

\begin{tcolorbox}[breakable,colback=teal!4,colframe=teal!35!black,title={Before you move on}]
What does \ml{അസ്വസ്ഥത} mean? (\textquotedblleft{}Unease; restlessness\textquotedblright{}.) Say it once more.
\end{tcolorbox}
\section[\texorpdfstring{\ml{വിതുമ്പുക} (vitumpuka)}{വിതുമ്പുക (vitumpuka)}]{\ml{വിതുമ്പുക} (vitumpuka) — to sob; to choke up}

\label{lesson:ML-C520-vitumpuka}



\begin{quote}
Before the new one: say the Malayalam for to throb, then the Malayalam for unease.
\end{quote}

\subsection*{You'll want to know: \ml{വിതുമ്പുക}}
\textbf{\ml{വിതുമ്പുക}} — \emph{vitumpuka} — \textquotedblleft{}to sob; to choke up\textquotedblright{}.

\begin{quote}
\ml{അച്ഛൻ} \ml{യാത്ര} \ml{പറഞ്ഞപ്പോൾ} \ml{കുട്ടി} \ml{വിതുമ്പി}.
\end{quote}

\begin{quote}
\emph{acchan yātra paṟaññappōḷ kuṭṭi vitumpi.}
\end{quote}

\begin{quote}
When Father said goodbye, the child began to sob.
\end{quote}

Past: \ml{വിതുമ്പി}, \emph{vitumpi}.

Crying held back, with trembling lips; louder crying is \ml{കരയുക}.

\begin{cousinweb}
\textbf{\ml{വിതുമ്പുക}} (\emph{vitumpuka}) is native Dravidian: DEDR 5399 gives the Malayalam verb with the sense of starting to cry, as a child does (DEDR 5399).
\end{cousinweb}

\begin{grammarlens}[title={What you've built}]
One more word for describing what you lived and felt.
\end{grammarlens}

\subsection*{Your turn}
\begin{itemize}
\raggedright
  \item \emph{Say it:} \emph{vitumpuka}
  \item \emph{Say it:} \emph{vitumpuka}, once more
  \item \emph{Say it:} \emph{acchan yātra paṟaññappōḷ kuṭṭi vitumpi.} — When Father said goodbye, the child began to sob.
  \item \emph{Recall:} say the Malayalam for to throb, then the Malayalam for unease, then say \emph{vitumpuka} again
\end{itemize}

\begin{tcolorbox}[breakable,colback=teal!4,colframe=teal!35!black,title={Before you move on}]
What does \ml{വിതുമ്പുക} mean? (\textquotedblleft{}To sob; to choke up\textquotedblright{}.) Say it once more.
\end{tcolorbox}
